Heterogeneous robot teams must form coalitions whose combined capabilities cover each task's requirements and
whose members arrive together before the task can start. This couples coalition formation, routing and scheduling,
and learned policies have recently been reported as the state of the art for it. We ask whether that claim survives a
comparison at equal compute. We present an adaptive large neighbourhood search (ALNS) over a deadlock-free plan
representation, a global task order with a minimal covering coalition per task, whose evaluator reproduces the
benchmark's simulator exactly and scores an insertion move in time linear in the coalition size. On the HeteroMRTA
benchmark, with up to 150 robots and 500 tasks, we compare it with the published reinforcement-learning policy,
three methods built on the constraint solver CP-SAT (including a fully parallel large neighbourhood search), the exact
MILP the benchmark used (CTAS-D, ported to an open solver and checked against the original authors' solutions), and
greedy restarts.
All methods run on the same pinned cores at the benchmark's own time budgets, and the confirmatory test run follows a
pre-registration frozen before it. ALNS is significantly better than every competitor on \ConeBheld{} of the 8
settings: its plans are \mbox{\ConeBclasslo--\ConeBclasshi\%} shorter than those of the CP-SAT-based methods and
greedy restarts, and \mbox{\ConeBrlgainlo--\ConeBrlgainhi\%} shorter than those of the learned policy, which runs at its
published time budget on the same CPUs. Making CP-SAT fully parallel does not close the gap, and the exact MILP finds a plan within the budget
for only \CtasSolved{} of the \CtasRuns{} instances with up to 200 tasks. On one core and in 2\,s, ALNS is already significantly better than every eight-core competitor on
\Ctwoheld{} of the 8 settings.
